\section{Implementación Física y Pruebas en Laboratorio}

Una vez validado el diseño lógico, se cargó el archivo de programación al chip FPGA de la placa de desarrollo. Las interconexiones en el circuito de pruebas se validaron empleando instrumentación básica de laboratorio.

A continuación se presentan las evidencias fotográficas de la implementación del hardware y de las señales medidas mediante osciloscopio y analizador lógico:

\placeholderfigure{Fotografía general del alambrado físico del sistema conectado a la FPGA y protoboard.}{imagenes/foto_alambrado_general.png}
\label{fig:foto_alambrado_general}

\placeholderfigure{Fotografía a detalle de la etapa de aislamiento de potencia (optoacoplador PC817 y relevador para bomba).}{imagenes/foto_etapa_potencia.png}
\label{fig:foto_etapa_potencia}

\placeholderfigure{Fotografía a detalle de las conexiones del convertidor ADC MCP3008 con la red de sensores.}{imagenes/foto_sensor_adc.png}
\label{fig:foto_sensor_adc}

\placeholderfigure{Medición en osciloscopio del ciclo de trabajo de la señal PWM generada para el servomotor (Posiciones de 1.0 ms y 2.0 ms).}{imagenes/medicion_osciloscopio_pwm.png}
\label{fig:medicion_osciloscopio_pwm}

\placeholderfigure{Captura del analizador lógico mostrando el protocolo SPI (lectura síncrona de los canales del ADC).}{imagenes/medicion_analizador_logico_spi.png}
\label{fig:medicion_analizador_logico_spi}
